\documentclass[a4paper]{article}
% UTF-8
\usepackage[T1]{fontenc}
\usepackage[utf8]{inputenc}

% Czcionka, która sensownie wygląda na monitorze.
\usepackage{lmodern}

% Język
\usepackage[polish]{babel}
% \usepackage[english]{babel}

% Rózne przydatne paczki:
% - znaczki matematyczne
\usepackage{amsmath, amsfonts}
% - wcięcie na początku pierwszego akapitu
\usepackage{indentfirst}
% - komenda \url 
\usepackage{hyperref}
% - dołączanie obrazków
\usepackage{graphics}
% - szersza strona
\usepackage[nofoot,hdivide={2cm,*,2cm},vdivide={2cm,*,2cm}]{geometry}
\frenchspacing
% - brak numerów stron
\pagestyle{empty}

% dane autora
\author{Krystian Czarnota}
\title{Przykładowy plik w systemie \LaTeX}
\date{\today}

% początek dokumentu
\begin{document}
\maketitle
To jest przykładowy akapit.

To jest nowy akapit.\\
To jest ten sam akapit po znaku\footnote{Znak nowej linii składa się z dwóch znaków.} nowej linii.

To jest fragment pracy naukowej:
\medskip

We say that a co-B\"uchi automaton ${\cal{A}}=\langle \Sigma , Q, q_0, \delta , \alpha\rangle$ is in the {\bf normal form} iff for each $\langle q, a, q'\rangle \in \delta$ if $q$ is in $\alpha$, then also $q'$ is in $\alpha$. Note that for a given NCW 
${\cal{A}}=\langle \Sigma, Q, q_0, \delta , \alpha\rangle$
the automaton ${\cal{A'}}=\langle \Sigma, Q', \langle q_0, 0\rangle, \delta', \alpha\times\{1\}\rangle$, where 
$Q'=Q\times\{0\} \cup \alpha \times \{1\}$ and 
$\delta'=\{\langle \langle q, i\rangle, a, \langle q', j \rangle\rangle \; | \; \langle q, a, q'\rangle \in \delta \wedge i \leq j 
\wedge  \langle q, i\rangle,\langle q', j \rangle\in Q' \}$ 
is in the normal form, recognizes the same language and has at most $2|Q|$ states.

 

\bigskip

 
 Na kolejnych stronach jest kawałek pracy ze strukturą.
\newpage
 
 \begin{abstract}
\maketitle TYTUŁ NAD TYTUŁAMI\\
michał, michał, michał, michał, co z ciebie wyrosnie??
\end{abstract}

\section{Wprowadzenie}
\subsection{Czas uniksowy}
Sposób zapisywania daty i czasu jako liczby sekund, które upłynęły od 1 stycznia 1970 roku o godzinie 00:00:00 UTC.
\subsection{Problem roku 2038}
Problem roku 2038 wynika ze sposobu przechowywania czasu uniksowego w niektórych starszych systemach komputerowych. W systemach wykorzystujących 32-bitową liczbę całkowitą ze znakiem do zapisywania liczby sekund maksymalna dodatnia wartość wynosi 2147483647.


\section{Indukcja matematyczna}\label{prelimeracja}
Chcemy pokazać, że dla każdego $n \in \mathbb{N}$ zachodzi $\sum_{i=0}^n 2^i + 1 = 2^{n+1}$. Zatem zdefiniujmy następujący zbiór:
\[
X = \{n \in \mathbb{N} : \sum_{i=0}^n 2^i + 1 = 2^{n+1}\}
\]

\subsection{Podstawa indukcji}
Dla $n = 0$
\[
\sum_{i=0}^0 2^i + 1 = 2^{0+1}
\]
\[
2^0 + 1 = 2^1
\]
\[
2 = 2
\]

zatem $0 \in X$

\subsection{Założenie indukcyjne}

Weźmy dowolne $n \in \mathbb{N}$ dla którego zachodzi $\sum_{i=0}^n 2^i + 1 = 2^{n+1}$\\
Pokażemy, że zachodzi dla $n+1$

\subsection{Krok indukcji}

Teza: $\sum_{i=0}^{n+1} 2^i + 1 = 2^{n+2}$

Skorzystajmy z założenia indukcyjnego.
\[
\sum_{i=0}^n 2^i + 1 = 2^{n+1}
\]
\[
\sum_{i=0}^n 2^i + 2^{n+1} + 1 = 2^{n+1} + 2^{n+1}
\]
\[
\sum_{i=0}^{n+1} 2^i + 1 = 2^{n+2}
\]

Teza zachodzi zatem $n+1 \in X$\\
z czego wynika, że $X=\mathbb{N}$

% Osadzona bibliografia; zwykle bibliografię się generuje narzędziem bibtex, ale jest ono poza ramami tego przedmiotu.
\begin{thebibliography}{12}
	\bibitem{ANB96}
	H. Andreka, I. Nemeti, J. van Benthem. Modal languages and bounded fragments of predicate logic. \emph{Journal of Philosophical Logic}, 27:217-274, 1998. 
	\bibitem{AFA97} 
	S. Arora, R. Fagin. On winning strategies in Ehrenfeucht--Fraiss'e games. \emph{Theoretical Computer Science}, 174(1--2):97--121, 1997. 
	\bibitem{BEN76} 
	J. van Benthem. Modal Correspondence Theory, \emph{dissertation, Mathematical Institute, University of Amsterdam}, 1976.
	\bibitem{BEN84} 
	J. van Benthem. 'Correspondence Theory'. In \emph{D. Gabbay, F. Guenthner, Handbook of Philosophical Logic, vol. II}, Reidel, Dordrecht, 1984.
	\bibitem{FLA77} 
	M. J. Fischer, R. E. Ladner. Propositional modal logic of programs. In \emph{Proceedings of the Ninth Annual ACM Symposium on theory of Computing (Boulder, Colorado, United States, May 04 - 04, 1977). STOC '77. ACM}, New York, NY, 286-294, 1977.
	\bibitem{GKV97}
	E. Gr\"{a}del, P. G. Kolaitis, and M. Y. Vardi. On the decision problem for two-variable first-order logic. \emph{Bulletin of Symbolic Logic}, 3(1):53--69, Mar, 1997.
	\bibitem{GRA98} 
	E. Gr\"{a}del. On the restraining power of guards. \emph{Journal of Symbolic Logic}, 64(4):1719-1742, 1999.
	\bibitem{KAZ06}
	Y. Kazakov. Saturation-Based Decision Procedures for Extensions of the Guarded Fragment.
	\emph{PhD thesis}, Universit\"{a}t des Saarlandes, Saarbr\"{u}cken, Germany, March 2006.
	\bibitem{PRA78}
	V. R. Pratt 1978. A practical decision method for propositional dynamic logic (Preliminary Report). In \emph{Proceedings of the Tenth Annual ACM Symposium on theory of Computing (San Diego, California, United States, May 01 - 03, 1978)}. STOC '78. ACM, New York, NY, 326-337, 1978.
	\bibitem{SCO62} 
	D. Scott. A decision method for validity of sentences in two variables. \emph{J. Symb. Logic} 27, 477, 1962.
	\bibitem{SZT04} 
	W. Szwast, L. Tendera. The guarded fragment with transitive guards. \emph{Annals of Pure and Applied Logic}, 128, 227-276, 2004.
\end{thebibliography}

\end{document}
